\section{Implementation}

\subsection{Project structure and runtime responsibilities}
The repository follows a deliberately flat layout. There is no framework, no bundler, and
no \texttt{node\_modules}.
\begin{lstlisting}[language=bash,numbers=none,caption={Relevant Rùa Toán Basic project layout.}]
ruatoan_basic/
|-- main.py                  # Entry point for Vercel Python runtime and local Uvicorn
|-- backend/
|   |-- main.py              # FastAPI app, API routes, security headers, lifespan hooks
|   |-- content.py           # JSON catalogue loader with mtime-keyed in-memory cache
|   |-- auth.py              # PBKDF2 hashing, HMAC session cookies, token handling
|   `-- database.py          # psycopg connection pool + schema initialization
|-- public/
|   |-- index.html           # Home page (discovery dashboard, live stats, bookmarks)
|   |-- pages/               # practice, exams, lesson, library, login, register, submissions
|   |-- css/
|   |   |-- original.css     # Authentic Rùa Toán layout, typography, Bungee badges, card grid
|   |   `-- style.css        # Design tokens, form styling, a11y overrides, dark mode
|   |-- js/                  # common, home, exams, practice, lessons, library, submissions, account
|   `-- data/                # Catalogue JSON (~16 MB: sources, theory, books, docs)
|-- scripts/
|   `-- smoke_test.py        # Automated end-to-end smoke test suite (all 8 flows)
|-- vercel.json              # Function includeFiles ("data/**") and maxDuration settings
`-- requirements.txt         # Pinned Python dependencies (fastapi, psycopg-pool, etc.)
\end{lstlisting}

\subsection{Content loading, indexing, and search}
At runtime, \texttt{backend/content.py} reads JSON files from \texttt{data/}. Parsed
values are cached by file modification time: a tuple of \texttt{(path, mtime)} is the cache
key, so a redeploy or a file change refreshes the cache automatically while repeated
requests reuse the parsed data. The loader reads \texttt{data/sources/*.json}
(excluding files whose name starts with \texttt{\_}), \texttt{data/theory/*.json},
\texttt{books.json}, and \texttt{docs.json}.

For search, helper \texttt{fold\_text} in \texttt{backend/main.py} normalises Vietnamese
diacritics (NFD decomposition, dropping combining marks, mapping \texttt{đ} to \texttt{d})
and lowercases. This is done \emph{once per problem at load time}: each problem gets a
\texttt{\_search} field containing the folded text of its label, statement, collection
title, and source. Request-time search is then a simple substring test against that field,
which keeps \texttt{/api/problems} fast even though it scans the in-memory list.

\texttt{/api/collections} returns summary objects (id, title, grade, year, kind, province,
problem count) without the problem arrays, so a page of 20 collections is small. The full
problem list of a collection is fetched only by \texttt{/api/collections/\{id\}}, which
supports \texttt{?page} and \texttt{?per} for pagination.

\subsection{Statement and mathematics rendering}
Content is stored as plain text. The frontend escapes it with \texttt{escapeText} (a
\texttt{<span>} with \texttt{textContent} assignment) before insertion, so no HTML from
the source can execute. Mathematical notation is written in LaTeX delimiters
(\texttt{\$...\$}, \texttt{\$\$...\$\$}) inside the plain text. \texttt{common.js}
configures MathJax 3 (loaded from the jsdelivr CDN) and exposes \texttt{renderMath()},
which waits for MathJax to be ready before calling \texttt{typesetPromise}. Every page
that displays mathematics calls \texttt{renderMath()} after inserting its content, so
formulas appear correctly even when the CDN is slow.

\subsection{Page composition and interactions}
\begin{description}
  \item[Home (\texttt{index.html} + \texttt{home.js}.)] Fetches \texttt{/api/stats} and
        renders the five totals; reads \texttt{savedProblems} from \texttt{localStorage}
        and renders the marked-problem list with a remove button per item.
  \item[Practice (\texttt{practice.html} + \texttt{practice.js}.)] A filter form
        (text, grade, topic) triggers \texttt{loadProblems(true)}, which calls
        \texttt{/api/problems}. Results are 20 per page; the ``Xem thêm'' button loads
        the next page. Each problem card has a mark/unmark toggle that writes to
        \texttt{localStorage} and updates its own label.
  \item[Collections (\texttt{exams.html} + \texttt{exams.js}.)] Lists collections 20 per
        page with a ``Xem thêm'' button. The detail view (\texttt{?id=...}) shows
        problems in order, a source link, and mark buttons.
  \item[Lessons (\texttt{lesson.html} + \texttt{lessons.js}.)] Loads all ninety lessons once,
        then filters client-side by title/summary. Each title is a toggle button that
        expands the lesson body and calls \texttt{renderMath()}.
  \item[Library (\texttt{library.html} + \texttt{library.js}.)] Loads books and documents
        in parallel with \texttt{Promise.allSettled}, so one failing endpoint does not
        blank the other. Client-side search over title/author/description.
  \item[Account (\texttt{login.html}/\texttt{register.html} + \texttt{account.js}.)] One
        form posts to \texttt{/api/login} or \texttt{/api/register} depending on the page;
        shows a success message and redirects after 700 ms.
  \item[Submissions (\texttt{submissions.html} + \texttt{submissions.js}.)] Lists the
        user's photos (with \texttt{loading="lazy"}), uploads 1--6 files with client-side
        size/count checks, and deletes with a confirmation dialog.
\end{description}

The shared \texttt{common.js} injects the navigation bar, checks \texttt{/api/me}, wires
logout, and provides helpers used by every page: \texttt{escapeText}, \texttt{safeUrl},
\texttt{api} (fetch with status-aware errors), \texttt{readSaved}/\texttt{writeSaved},
\texttt{showMessage}, and \texttt{renderMath}.

\subsection{Authentication}
\texttt{backend/auth.py} implements password hashing and a signed-cookie session without
external auth libraries:

\begin{itemize}
  \item \texttt{hash\_password} relies on \texttt{hashlib.pbkdf2\_hmac} with SHA-256 and
        200,000 iterations, using a 16-byte \texttt{os.urandom} salt, and stores
        \texttt{pbkdf2\_sha256\$200000\$<salt>\$<hash>} (hex).
  \item \texttt{check\_password} accepts both the new prefixed format and the legacy bare
        hex format, so old hashes keep working if the cost is raised later.
  \item \texttt{make\_cookie} and \texttt{read\_user\_id} sign a structured body
        (\texttt{<user\_id>:<expiry>}) with HMAC-SHA256 using \texttt{SESSION\_SECRET}; the
        cookie is \texttt{HttpOnly}, \texttt{SameSite=Lax}, \texttt{Secure} on HTTPS, and
        lives 14 days.
\end{itemize}

\texttt{/api/register} validates name (1--80 chars), email (email-validator), and password
(8--128 chars), checks uniqueness case-insensitively, and inserts the user.
\texttt{/api/login} returns a generic ``Email hoặc mật khẩu chưa đúng'' on failure (no user
enumeration) and sets the session cookie on success. \texttt{/api/logout} clears the cookie.
\texttt{/api/me} returns the current user or \texttt{null}.
\texttt{/api/change-password} verifies current credentials and updates the password hash.
For account recovery, \texttt{/api/request-reset} generates a cryptographically random
32-byte token and stores its SHA-256 hash in \texttt{password\_resets} with a 30-minute
expiration interval. \texttt{/api/reset-password} verifies the single-use token against
its database record, updates the user's password hash, and purges the reset entry upon
successful completion.

\subsection{Database operations and photo API}
\texttt{backend/database.py} keeps a module-level \texttt{psycopg\_pool.ConnectionPool}
(min 1, max 5) created lazily. Strict connection pooling is enforced to prevent
connection starvation under serverless workloads. Schema creation (\texttt{prepare\_database})
runs once in the FastAPI \texttt{lifespan} hook, not on every request.

The photo endpoints:

\begin{itemize}
  \item \texttt{GET /api/photos} returns the logged-in user's photos (newest first) with
        metadata only.
  \item \texttt{POST /api/photos} accepts a multipart form with \texttt{files} (1--6),
        validates count, per-file size (1.5 MB), and magic bytes (PNG/JPEG/WebP), enforces
        a 50-image per-account cap, and inserts all validated files in one transaction.
  \item \texttt{GET /api/photos/\{id\}} streams the stored bytes with
        \texttt{X-Content-Type-Options: nosniff}.
  \item \texttt{DELETE /api/photos/\{id\}} deletes only a row owned by the current user.
\end{itemize}

\subsection{Deployment and environment}
Deployment is managed on Vercel using the Python runtime. Root \texttt{main.py} serves as the
application entry point, exposing the FastAPI ASGI \texttt{app} instance. The serverless
configuration in \texttt{vercel.json} specifies:
\begin{lstlisting}[numbers=none]
{
  "functions": {"main.py": {"includeFiles": "data/**", "maxDuration": 30}}
}
\end{lstlisting}
This configuration ensures that the 16 MB JSON catalogue under \texttt{data/} is
bundled into the serverless function execution sandbox, while static assets under \texttt{public/}
(HTML pages, CSS files, JavaScript bundles, images) are served directly by Vercel's global edge CDN
with zero compute overhead and no cold-start latency. Dynamic requests and \texttt{/api/*}
endpoints are handled natively by FastAPI.

Environment variables required in production: \texttt{DATABASE\_URL} and
\texttt{SESSION\_SECRET}. \texttt{.env.local} is gitignored and never committed.

\subsection{Bug fixes and hardening (improvement pass)}\label{sec:cleanup}
The first version of the basic site had many defects. The improvement pass addressed them
in four groups:

\begin{enumerate}
  \item \textbf{Frontend robustness.} Every \texttt{fetch} now checks \texttt{response.ok}
        and is wrapped in \texttt{try/catch}; failures render a message instead of leaving
        the page in a silent loading state. MathJax rendering waits for the library to load.
        \texttt{localStorage} access is guarded. Double-click protection was added to
        ``load more'' buttons. Photo delete asks for confirmation. Upload validates
        count and size before sending.
  \item \textbf{Security.} See Section~\ref{sec:security} (the list of headers, bounds,
        and magic-byte checks).
  \item \textbf{Backend correctness.} Connection pooling replaces one connection per
        request; schema creation runs once at startup; SQL uses \texttt{psycopg.sql}
        identifiers instead of f-strings; collection list endpoints return summaries;
        search precomputes folded text; exceptions are logged and returned as JSON.
  \item \textbf{Bundle size.} Redundant data copies under \texttt{public/} and a 2.5 MB
        backup file were removed, and two 1.2 MB screenshots were moved out of
        \texttt{public/}. The deployable data set dropped from 35 MB to 16 MB.
\end{enumerate}

\subsection{CSS architecture and interface unification}
The styling follows a unified, two-tier architecture designed to prevent style conflicts while
preserving the complete visual character of the original platform:
\begin{enumerate}
  \item \textbf{Core brand stylesheet (\texttt{original.css}):} Contains the structural layout
        definitions for the top header, hero banners (\texttt{phead}), stat counters, topic
        pills, collection card grids (\texttt{ccard}), pagination, and footer. It embeds the
        retro typography rules for \textbf{Bungee} and mathematical layout styles for KaTeX/MathJax.
  \item \textbf{Design system and component overrides (\texttt{style.css}):} Declares CSS
        custom properties in \texttt{:root} for semantic colors, spacing scales (\texttt{--s-1} to
        \texttt{--s-8}), border radiuses, and tactile 3D shadows. It styles components unique to
        the basic version (authentication forms, photo upload drop zones, image thumbnail grids,
        interactive message banners, bookmark trays) and enforces accessibility rules:
        high-contrast input backgrounds (\texttt{\#163523}) with light text (\texttt{\#eaf6e3}),
        prominent \texttt{:focus-visible} neon outlines (\texttt{\#b6ff3b}), and readability
        overrides for mathematics statement cards.
\end{enumerate}
All HTML pages include \texttt{style.css} followed by \texttt{original.css}, ensuring
predictable cascade order, zero visual conflicts, and full responsiveness across mobile (390 px)
and desktop viewports.

\subsection{Content preparation}
The content directory is reused unchanged from the original project; the crawl and
preparation scripts (Python/Node, Gemini/Claude extraction) are not part
of this version and is not run during page requests. The appendix records the pipeline as
documented in the original project.
